\documentclass[12pt]{article}
\usepackage[utf8]{inputenc}
\usepackage[french]{babel}
\usepackage{amsmath,amssymb}
\usepackage{xcolor}

\begin{document}

% slide borne supérieure

\begin{itemize}
    \renewcommand{\labelitemi}{\textbullet}
	\setlength\itemsep{1em}
    \item Le nombre de graduations $m(\ell)$ est minimisé lorsque la taille du\newline premier bloc est proche de $\sqrt{\ell}$
    \item $\sqrt{\ell}$ graduations pour le premier bloc, et $\sqrt{\ell}$ graduations éparpillées
    \item[] $\implies$ Notre solution est d'environ $\sqrt{\ell} + \sqrt{\ell} = 2\sqrt{\ell}$ graduations
\end{itemize}

\vspace{1cm}

% slide borne inférieure

Il y a au moins $\sqrt{2\ell}$ graduations, car pour $\sqrt{2\ell}$ graduations, on a

$\frac{\sqrt{2\ell}(\sqrt{2\ell}-1)}{2}$ distances possibles.

\vspace{1cm}

% slide victor

\begin{itemize}
    \renewcommand{\labelitemi}{\textbullet}
	\setlength\itemsep{1em}
    \item Objectif: trouver une fonction $f(\ell)=a\sqrt{\ell}+b$ qui approxime $m(\ell)$
    \item Soit $x=\sqrt{\ell}$, on cherche à approximer $m(x^2)$ par une fonction affine $ax+b$
    \item $f(\ell) = 1.721\sqrt{\ell}+0.208$ avec $R^2=0.998$ par régression linéaire
\end{itemize}

\vspace{1cm}

% slide solutions parfaites 1

\begin{itemize}
    \renewcommand{\labelitemi}{\textbullet}
	\setlength\itemsep{1em}
    \item $n = m(\ell) \implies \ell = \frac{n(n-1)}{2}$
    \item $d_i$ : distance entre les graduations $i$ et $i+1$
    \item $d_1 + d_2 + \cdots + d_{n-1} = \ell$ et tous les $d_i$ sont distincts
    \item[] $\implies$ $d$ est une permutation de $[1,\dots,n-1]$
\end{itemize}

\[
\hspace*{-1.5cm}
\color{blue!60!black}{d_1 + \cdots + d_{n-1}}
\;\ge\;
\color{green!40!black}{1 + 2 + \cdots + (n-1)}
\]

\vspace{-0.8cm}

\[
\hspace*{-1.5cm}
\color{green!40!black}{1 + 2 + \cdots + (n-1)} = \frac{n(n-1)}{2} = \ell
\]

\vspace{1cm}

% slide solutions parfaites 2

\begin{itemize}
    \renewcommand{\labelitemi}{\textbullet}
	\setlength\itemsep{1em}
    \item Objectif: ordonner le tableau $[1,\dots,n-1]$ pour trouver la solution parfaite
    \item $d_i = 1$ ne peut être qu'à côté de $d_j=n-1$ $\implies$ cette valeur se trouve au bord
    \item On a donc $d=[1, n-1, \ldots]$
	\item $d_i=2$ ne peut être qu'à côté de $d_j=n-1$ $\implies$ impossible
\end{itemize}

\end{document}
